\section{Phase-type state spaces}\label{sec:phasetype}

\paragraph{A dictionary of Erlang densities.} Fix a set of rates $\beta_1,\dots,\beta_K>0$ on a
logarithmic grid, so that each rate corresponds to a time scale $1/\beta_k$, and an Erlang order
$R\ge1$. For $k=1,\dots,K$ and $r=0,\dots,R-1$ we define the basis densities
\begin{equation}\label{eq:basis}
g_{kr}(t)=\beta_k\,\frac{(\beta_k t)^r}{r!}\,e^{-\beta_k t},\qquad t\ge0 .
\end{equation}
Each $g_{kr}$ is the density of a gamma (Erlang) distribution with shape $r+1$ and rate
$\beta_k$, so $\int_0^\infty g_{kr}(t)\,\mathrm{d}t=1$ and its mean is $(r+1)/\beta_k$. A kernel is a
nonnegative combination
\begin{equation}\label{eq:kernel}
\phi(t)=\sum_{k=1}^{K}\sum_{r=0}^{R-1}a_{kr}\,g_{kr}(t),\qquad a_{kr}\ge0 .
\end{equation}
Because each basis density integrates to one, the total mass of the kernel, which is its
branching ratio, is simply $\sum_{k,r}a_{kr}$. With $R=1$ the dictionary contains only
exponentials, and mixtures of exponentials can approximate any completely monotone kernel,
power laws included (Bernstein's theorem). Orders $r\ge1$ add kernels that start at zero and peak
later, which describe delayed reactions. More generally, phase-type densities are dense in the set
of densities on $\R_+$.

\paragraph{Exact recursion.} Every state in this paper is built from such Erlang chains: the
Hawkes excitation, the response to news, the neural latent state of EPT and the renewal channel of
EPT-X. For all of them, the evolution between two events and its integral are available in closed
form. The following lemma gives the formulas for one source and one rate.

\begin{lemma}[Erlang-chain recursion]\label{lem:recursion}
Let $t_1<t_2<\dots$ be the event times of one source, let $\beta>0$ be one rate of the
dictionary, and define the chain of states
\[
s_r(t)=\sum_{n:\,t_n<t} g_r(t-t_n),\qquad r=0,\dots,R-1,
\]
where $g_r$ is the basis density~(\ref{eq:basis}) with rate $\beta$. Write
$\bm s=(s_0,\dots,s_{R-1})^\top$, let $I$ be the $R\times R$ identity and let $B$ be the
$R\times R$ down-shift matrix ($B_{r,r-1}=1$ and all other entries zero).
\begin{enumerate}
\item Between events, $\dot{\bm s}=\beta(B-I)\,\bm s$; at an event, $s_0$ increases by $\beta$.
\item For a gap $\Delta\ge0$ with no event in $(t,t+\Delta)$, let $y=\beta\Delta$ and
$w_j=e^{-y}y^j/j!$. Then
\[
s_r(t+\Delta)=\sum_{j=0}^{r}w_j\,s_{r-j}(t),\qquad
\int_t^{t+\Delta}s_r(u)\,\mathrm{d}u=\frac{1}{\beta}\sum_{j=0}^{r}P(j+1,y)\,s_{r-j}(t),
\]
where $P(a,y)$ is the regularised lower incomplete gamma function.
\end{enumerate}
\end{lemma}

\begin{proof}
\emph{Step 1 (dynamics).} Differentiating~(\ref{eq:basis}) gives
\[
\frac{\mathrm{d}}{\mathrm{d}t}g_0(t)=-\beta\,g_0(t),\qquad
\frac{\mathrm{d}}{\mathrm{d}t}g_r(t)=\beta\bigl(g_{r-1}(t)-g_r(t)\bigr)\quad(r\ge1).
\]
Summing over past events gives $\dot s_0=-\beta s_0$ and $\dot s_r=\beta(s_{r-1}-s_r)$, which is
$\dot{\bm s}=\beta(B-I)\bm s$. Since $g_r(0)=\beta\,\mathbb{1}[r=0]$, a new event adds $\beta$ to
$s_0$ and leaves the other states unchanged.

\emph{Step 2 (evolution over a gap).} The solution is $\bm s(t+\Delta)=e^{\beta(B-I)\Delta}\bm s(t)$.
The matrices $B$ and $I$ commute and $B^R=0$, so
\[
e^{\beta(B-I)\Delta}=e^{-y}\,e^{yB}=e^{-y}\sum_{j=0}^{R-1}\frac{y^j}{j!}B^j .
\]
Because $B^j$ shifts a vector down by $j$ places, component $r$ of $e^{\beta(B-I)\Delta}\bm s(t)$
is $\sum_{j\le r}w_j s_{r-j}(t)$.

\emph{Step 3 (integral over a gap).} Integrating the weights over the gap,
\[
\int_0^{\Delta}e^{-\beta u}\frac{(\beta u)^j}{j!}\,\mathrm{d}u=\frac{1}{\beta}\,P(j+1,y),
\]
which gives the second formula.
\end{proof}

Evaluating a likelihood therefore costs $O(N d K R^2)$ operations for $N$ events in $d$
dimensions, and it involves no Monte Carlo or quadrature error.
